\documentclass[11pt,a4paper]{article}

% Sprache und Schriften (mit XeLaTeX kompilieren)
\usepackage{fontspec}
\usepackage[ngerman]{babel}

% Mathematik, Abbildungen und Links
\usepackage{amsmath, amssymb}
\usepackage{graphicx}
\usepackage[hidelinks]{hyperref}

\title{Titel des Artikels}
\author{Vorname Nachname}
\date{\today}

\begin{document}

\maketitle

\begin{abstract}
  Kurze Zusammenfassung der Arbeit: das Problem, der Ansatz und die wichtigsten Ergebnisse.
\end{abstract}

\section{Einleitung}
\label{sec:einleitung}

Dies ist ein Beispielartikel. Umlaute und das ß (ä, ö, ü, ß) funktionieren dank
XeLaTeX direkt. Mehr dazu in Abschnitt~\ref{sec:hauptteil} und in
Quelle~\cite{knuth1984}.

\section{Hauptteil}
\label{sec:hauptteil}

Eine nummerierte Gleichung:
\begin{equation}
  \label{eq:euler}
  e^{i\pi} + 1 = 0
\end{equation}

und ein ausgerichtetes System:
\begin{align}
  \nabla \cdot \vec{E} &= \frac{\rho}{\varepsilon_0}, \\
  \nabla \cdot \vec{B} &= 0.
\end{align}

Gleichung~\eqref{eq:euler} verbindet fünf grundlegende Konstanten.

\section{Fazit}

Fassen Sie hier die wichtigsten Ergebnisse zusammen.

\bibliographystyle{plain}
\bibliography{literatur}

\end{document}
